\section{Strategy of the proof}\label{sec:strategy}

The proof of \cref{thm:main} has two layers. The first is a short reduction to the boundedness of the
critical connection operator for graphs with a unimodular automorphism group
(\cref{sec:reduction}). The second proves this boundedness; it occupies
\cref{sec:prelim,sec:bridges,sec:maps,sec:moments,sec:triangle,sec:deformation,sec:escapes}. We first
explain the mechanism (\cref{sec:mechanism}). We then state the four steps of the proof precisely: two a
priori estimates on critical and near-critical clusters (Steps~1 and~2, \cref{sec:part1}), and two
properties of a deformation of the connection kernel (Steps~3 and~4, \cref{sec:part2}). Finally we show
how the four steps combine (\cref{sec:assembly}). \Cref{fig:dependencies} summarises the logical
structure.

\begin{figure}[t]
\centering
% Logical dependencies of the proof (a graph of labelled statements; nothing quantitative is drawn)
\begin{tikzpicture}[
    box/.style={draw, rounded corners=2pt, align=center, font=\small, text width=4.6cm,
                inner sep=3pt, minimum height=1.05cm,
                execute at begin node={\hyphenpenalty=10000\exhyphenpenalty=10000}},
    main/.style={box, very thick},
    cite/.style={box, draw=figgrey, dashed},
    arr/.style={-{Stealth[length=2.2mm]}, thick},
    lab/.style={font=\footnotesize\itshape, fill=white, inner sep=1pt}]
  % column headers
  \node[font=\footnotesize\scshape, text=figgrey] at (5.1,1.0) {Part I: a priori estimates};
  \node[font=\footnotesize\scshape, text=figgrey] at (13.05,1.0) {Part II: deformation};
  % Part I
  \node[box] (bridges) at (2.45,0)    {Bridges in critical clusters\\ \cref{prop:good}, \cref{sec:bridges}};
  \node[box] (ampl)    at (7.75,0)    {Heavy tail $\Rightarrow$ efficient maps\\ \cref{prop:amplification}, \cref{sec:maps}};
  \node[box] (eff)     at (2.45,-2.0) {Efficient maps $\Rightarrow$ light tail\\ \cref{prop:efficient}, \cref{sec:moments}};
  \node[box] (mom)     at (7.75,-2.0) {Step~1: \mbox{$\EE_{p_c}\abs{C_o}^\alpha<\infty$} for \mbox{$\alpha<\tfrac12$}\\ \cref{thm:moments}};
  \node[box] (apr)     at (7.75,-4.0) {Step~2: \mbox{$\chi(p)\le C_\eta\delta^{-1-\eta}$}, \mbox{$\nabla_p\le C_\eta\delta^{-\eta}$}\\ \cref{thm:apriori}};
  % Part II
  \node[box] (def)     at (13.05,0)    {Step~3: corridor deformation\\ \cref{prop:deformation}, \cref{sec:deformation}};
  \node[box] (esc)     at (13.05,-2.0) {Step~4: escapes at both ends\\ \cref{prop:escapes}, \cref{sec:escapes}};
  % assembly, the sprinkling bound and the main results
  \node[box]  (spr)    at (2.45,-6.0)  {Sprinkling bound\\ \cref{lem:sprinkling}};
  \node[main] (cb)     at (7.75,-6.0)  {Assembly: \mbox{$\norm{T_{p_c}}<\infty$} if \cref{ass:U} holds\\ \cref{thm:crit-bdd}, \cref{sec:assembly}};
  \node[cite] (nonu)   at (2.45,-8.0)  {Nonunimodular case\\ \cite[Theorem~2.9]{Hut19}};
  \node[box]  (aux)    at (13.05,-8.0) {\mbox{$\ptwo\le p_u$}; simultaneous interval\\ \cref{lem:p22-pu}, \cref{prop:interval}};
  \node[main] (thm)    at (7.75,-8.0)  {Main result\\ \cref{thm:main}};
  \node[cite] (hut)    at (13.05,-9.9) {Conditional results\\ \cite{DCT16,Hut20L2,Hut22sup}};
  \node[box]  (cons)   at (7.75,-9.9)  {Mean-field behaviour\\ \cref{thm:consequences}, \cref{sec:consequences}};
  % arrows inside Part I
  \draw[arr] (bridges) -- (ampl);
  \draw[arr] (ampl) -- node[lab, right=3pt] {by contradiction} (mom);
  \draw[arr] (eff) -- (mom);
  \draw[arr] (mom) -- (apr);
  \draw[arr] (apr) -- (cb);
  % Part II feeds the assembly through one bus
  \draw[thick] (def.west) -| (10.4,-6.0);
  \draw[thick] (esc.west) -- (10.4,-2.0);
  \draw[arr] (10.4,-6.0) -- (cb.east);
  % the sprinkling bound enters the assembly and the reduction
  \draw[arr] (spr) -- (cb);
  \draw[arr] (spr.south east) -- (thm.north west);
  % from critical boundedness to the main results
  \draw[arr] (cb) -- (thm);
  \draw[arr] (nonu) -- (thm);
  \draw[arr] (aux) -- (thm);
  \draw[arr] (thm) -- (cons);
  \draw[arr] (hut) -- (cons);
\end{tikzpicture}

\caption{Logical structure of the proof. An arrow points from a statement to a statement whose proof
uses it. Dashed boxes are results from the literature. Step~1 is proved by contradiction: if the
critical cluster size had an infinite moment of order $\alpha<1/2$, the amplification would produce
efficient endpoint maps, which force all such moments to be finite. The assembly
(\cref{sec:assembly}) combines the four steps with the sprinkling bound.}\label{fig:dependencies}
\end{figure}

\subsection{Reduction to critical boundedness}\label{sec:reduction}

The core of the argument is the following theorem.

\begin{theorem}[Critical boundedness]\label{thm:crit-bdd}
Assume \cref{ass:G,ass:U}. Then $\norm{T_{p_c}}<\infty$.
\end{theorem}

Everything else in \cref{thm:main} follows from \cref{thm:crit-bdd} and the tools of
\cref{sec:prelim}.

\begin{proof}[Proof of \cref{thm:main} given \cref{thm:crit-bdd}]
We first show that $T_{p_c}$ is bounded and $p_c<\ptwo$. If $\Aut(G)$ is unimodular,
$\norm{T_{p_c}}<\infty$ by \cref{thm:crit-bdd}. Since $p_c<1$ (\cref{lem:pc}), the sprinkling bound
(\cref{lem:sprinkling}) with a small $a>0$ then shows that $T_p$ is bounded at
$p=p_c+(1-p_c)a>p_c$, so $p_c<\ptwo$. If $\Aut(G)$ is not unimodular, then $p_c<\ptwo$ by a theorem of
Hutchcroft (\cref{prop:nonunimodular}). Since $T_p\le T_{p'}$ entrywise for $p\le p'$, the kernel
$T_p$ is bounded for every $p<\ptwo$ (\cref{lem:kernels}), in particular for $p=p_c$.

\Cref{lem:p22-pu} gives $\ptwo\le p_u$, and \cref{lem:kernels} identifies $\norm{T_{p_c}}$ with
$\sup_{p<p_c}\norm{T_p}$. This proves~\eqref{eq:main}. Finally, choose $p_c<p_1<p_2<\ptwo$. Then $T_{p_2}$
is bounded, and \cref{prop:interval} yields the simultaneous statement on $[p_1,p_2]$.
\end{proof}

\subsection{The mechanism: row sums versus norm}\label{sec:mechanism}

\paragraph{Row sums and norm.}
The connection kernel $T_p$ is nonnegative, symmetric and invariant, and its row sums equal $\chi(p)$.
By the Schur test (\cref{lem:kernels}), $\norm{T_p}\le\chi(p)$. Below $p_c$ the susceptibility is finite
by sharpness of the phase transition, so $T_p$ is bounded. The row sums, however, explode as
$p\uparrow p_c$: if $\chi$ stayed bounded, then $\chi(p_c)<\infty$ by monotone convergence, and the
sprinkling bound (\cref{lem:sprinkling}) would give $\chi(p)<\infty$ for some $p>p_c$, which is
impossible because $\theta(p)>0$. \Cref{thm:crit-bdd} therefore asserts that the norm of $T_p$ stays
bounded although its row sums do not. On a nonamenable graph this is possible because the $\ell^2$ norm
can be much smaller than the row sums. The basic example is the \emph{lazy simple random walk}
\begin{equation*}
  P(x,y)\coloneqq\frac{1}{2d}\,\1\{x\sim y\}+\frac12\,\1\{x=y\}\,,
  \qquad
  \rho\coloneqq\norm{P}<1\,,
\end{equation*}
whose row sums are one, while its norm is strictly smaller than one by nonamenability
(\cref{lem:walk}). We use $P$ and $\rho$ throughout.

\paragraph{What unboundedness would mean.}
Suppose that $\norm{T_{p_c}}=\infty$, and write $\delta\coloneqq p_c-p$ for $p<p_c$. Adding an
independent configuration of parameter $a$ to one of parameter $p$ and splitting open paths at the added
edges gives the sprinkling bound (\cref{lem:sprinkling})
\begin{align*}
  T_{p+(1-p)a}&\le\sum_{j\ge0}a^j\,T_p(\Adj\,T_p)^j\qquad\text{entrywise, and hence}\\
  \norm{T_{p+(1-p)a}}&\le\frac{\norm{T_p}}{1-ad\norm{T_p}}\qquad\text{if }ad\norm{T_p}<1\,.
\end{align*}
With $a=\delta/(1-p)$, so that $p+(1-p)a=p_c$, this turns a small norm at~$p$ into boundedness
at~$p_c$. Hence
\begin{equation}\label{eq:norm-lower}
  \norm{T_p}\ge\frac{1-p}{d\,\delta}\qquad\text{for all }p<p_c\,.
\end{equation}
Step~2 below shows that $\chi(p)\le C_\eta\delta^{-1-\eta}$ for every $\eta>0$. Hence the ratio
\begin{equation}\label{eq:R}
  R=R(p)\coloneqq\frac{\chi(p)}{\norm{T_p}}\in[1,\infty)
\end{equation}
satisfies $R\le C_\eta\delta^{-\eta}$ for every $\eta>0$: the norm of $T_p$ would agree with its row
sums up to a factor that is smaller than every power of $1/\delta$. We show that this is impossible by
deforming $T_p$ in a way that decreases its row sums by a large factor while decreasing its norm by at
most a factor two. This contradicts $\norm{\,\cdot\,}\le$ row sum.

\paragraph{Three structural facts.}
The comparison of norm and row sums rests on three properties of unimodular transitive graphs and of the
connection kernel, which we use throughout. First, the \emph{mass-transport principle}
(\cref{lem:mtp}): $\sum_yF(o,y)=\sum_xF(x,o)$ for every nonnegative function $F$ on $V\times V$ that is
invariant under the diagonal action of $\Aut(G)$; in particular every invariant kernel has equal row and
column sums. Second, for $p<p_c$ the kernel $T_p$ is \emph{positive semidefinite}
(\cref{lem:subcritical}). Third, the diagonal entry $\tr S\coloneqq S(o,o)$ of a bounded invariant
operator is a \emph{trace}: it is cyclic, and $\abs{\tr(BA)}\le\norm B\tr A$ for $A$ positive
semidefinite and $B$ self-adjoint (\cref{lem:trace}). Two consequences show how the decay of the
random walk on $\ell^2(V)$ enters. By the Cauchy--Schwarz inequality for $T_p$ and the trace,
\begin{equation*}
  (P^mT_p)(a,b)^2\le(P^mT_pP^m)(a,a)\,T_p(b,b)=\tr(P^{2m}T_p)\le\rho^{2m}\,:
\end{equation*}
a random walk of $m$ steps smooths the connection kernel at rate $\rho^m$, however large $\norm{T_p}$ is
(\cref{lem:smoothing}). And sums of the form $\sum_vP^k(o,v)\,T_p^3(v,o)=\tr(P^kT_p^3)$ are at most
$\rho^k\tr(T_p^3)=\rho^k\nabla_p$; this is how the triangle diagram enters the proof.

\paragraph{A deformation that separates the two.}
Fix an integer $k\ge3$ and let $w$ be a path of $k$ steps of the lazy walk, with probability $\mu(w)$
under~$P$. Erasing its loops produces a self-avoiding path $\gamma_w$ from $u=w_0$ to $v=w_k$. On the
\emph{corridor event} $\Cor_w$ all edges of $\gamma_w$ are open and all other edges at the interior
vertices of $\gamma_w$ are closed. Then $\gamma_w$ is an isolated open path, a \emph{corridor}. It joins
the clusters of $u$ and $v$ outside its interior, and it is their only open link unless these two
clusters coincide (\cref{fig:corridor}).
The deformation multiplies the indicator of $\{x\conn y\}$ by $\exp(-t\,c(w))$ for every corridor $w$
all of whose edges are pivotal for $x\conn y$, where $t\ge0$ is the deformation time and the weight
$c(w)=\mu(w)/\PP_p(\Cor_w)$ is calibrated so that each corridor is penalised at rate equal to its
random-walk probability. A penalised connection consists of a connection from $x$ to~$u$, a corridor,
and a connection from $v$ to~$y$. Summed over corridors, the rate at which the deformed kernel $T_t$
decreases is therefore at most $2\,T_pP^kT_p$, whose norm is at most $2\rho^k\norm{T_p}^2$. The rate
at which its row sum decreases is, roughly, the $\mu$-average of the product of the cluster masses at
$u$ and~$v$. Here no factor comparable to $\rho^k$ appears, since the row sums of $P^k$ equal one. Thus
over a time interval of length $(kR)^{C}/\norm{T_p}$ the norm changes by a negligible amount, while
the row sum, if the masses at the ends of corridors are comparable to the full row sum, decreases at a
rate proportional to its own square and must drop far below $\norm{T_p}$.

\paragraph{The three costs.}
Turning the last sentence into a proof requires three corrections, and each has a price.
\begin{enumerate}[label=(\roman*)]
  \item \emph{Distinct clusters.} The edges of a corridor are pivotal for $x\conn y$ only if the
  clusters of $u$ and $v$ outside the corridor are distinct. By the van den Berg--Kesten inequality
  and the trace bound above, the resulting error is at most
  $C\chi(p)^2\rho^k\nabla_p=CR^2\nabla_p\rho^k\norm{T_p}^2$
  (\cref{prop:deformation}\ref{it:def-rate}). It is negligible once $k$ is a large multiple of
  $\log R+\log\nabla_p$.
  \item \emph{Removing the corridor.} The cluster masses at $u$ and $v$ are computed after deleting the
  interior of the corridor, which may cut off most of the mass. Step~4 shows that for a positive
  fraction of paths $w$, both endpoints are joined by short paths avoiding the corridor to vertices
  whose connections to the neighbourhood of the corridor are tiny. Its key input is the smoothing bound
  above. Transporting the mass along these short paths costs a factor of order $(kR)^{-C}$.
  \item \emph{Interaction of corridors.} Penalties of different corridors on the same connection
  interact, which costs a factor $\e^{-tc_k}$ with $c_k=\e^{C_0k}$. This is harmless as long as
  $tc_k\ll1$ up to the final time $t\approx(kR)^C/\norm{T_p}$. By~\eqref{eq:norm-lower} this holds if
  $\e^{C_0k}(kR)^C\ll1/\delta$, in particular if $k$ and $\log R$ are $o(\log(1/\delta))$.
\end{enumerate}
The requirements (i) and (iii) can be met simultaneously if $\log R$ and $\log\nabla_p$ are
$o(\log(1/\delta))$: then $k$ can be chosen as a large multiple of
$\log R+\log\nabla_p+\log\log(1/\delta)$, which is still $o(\log(1/\delta))$. A bound
$R+\nabla_p\le C\delta^{-a}$ with a fixed $a>0$ would also do, but only if $a$ is smaller than a
threshold determined by constants of the construction (such as $C_0$ and the constants of Step~4) that
we do not control explicitly. This is why Step~2 proves the a priori bounds with an arbitrarily small
power of $1/\delta$.

\paragraph{Where the a priori bounds come from.}
Mean-field theory predicts $\PP_{p_c}(\abs{C_o}\ge n)\asymp n^{-1/2}$; \cref{thm:consequences} confirms
this a posteriori. Step~1 proves the weaker statement that $\EE_{p_c}\abs{C_o}^\alpha<\infty$ for every
$\alpha<1/2$. Step~2 converts it into the near-critical bounds. A decision-tree inequality transfers
the critical tail to $\chi(p)\le\delta^{-1-o(1)}$. A differential inequality for the triangle diagram
gives $\nabla_p\le\delta^{-o(1)}$; in it, the critical tail bounds the probability that two adjacent
distinct clusters are both large, and a concentration estimate bounds the number of edges between
them.

\paragraph{Critical moments by contradiction.}
Step~1 is proved by contradiction. If $\EE_{p_c}\abs{C_o}^\alpha=\infty$ for some $\alpha<1/2$, then at
infinitely many scales $n$ the probability that $\abs{C_o}\in[n,2n)$ is at least $n^{-\beta}$ for some
$\beta<1/2$: there are more large critical clusters than mean-field theory allows. We exploit them by
splicing independent small critical clusters onto a large one. This produces \emph{local endpoint
maps}, random rules that select a vertex of the cluster of the root, whose selection kernel has a very
small norm on $\ell^2(V)$ compared with their cost, measured by the logarithms of the inverse success
probability and of the cluster size. The multiplicity of the splicing is controlled by a bound on the
number of bridges in pieces cut out of critical clusters. Running such an efficient map slightly below
$p_c$ shows that $\chi$ grows at least as fast as a solution of $\chi'=\chi^{2-\gamma}$ for every
$\gamma>0$. This gives $\chi(p)\le\delta^{-1-o(1)}$ and hence a near-mean-field tail bound for the
critical cluster size, which contradicts the assumption. The susceptibility bound obtained in this way
is conditional on the heavy tail, which is used to construct the maps; Step~2 proves it unconditionally
from Step~1.

\subsection{Part I: a priori estimates (Steps 1 and 2)}\label{sec:part1}

\begin{theorem}[Step~1: critical moments]\label{thm:moments}
Assume \cref{ass:G,ass:U}. Then $\EE_{p_c}\abs{C_o}^\alpha<\infty$ for every $\alpha\in(0,1/2)$.
\end{theorem}

The proof of \cref{thm:moments} combines three propositions. The first concerns the geometry of
critical clusters. A \emph{bridge} of a finite connected graph~$H$ is an edge whose removal
disconnects~$H$. For vertices $x,y$ of~$H$ let $\dbr^H(x,y)$ be the number of bridges of~$H$ that
separate $x$ from~$y$; this is a pseudometric, the \emph{bridge distance}, and
$\diambr(H)\coloneqq\max_{x,y}\dbr^H(x,y)$ is the \emph{bridge diameter}. Throughout, a vertex set $B$
of an open cluster is regarded as a graph with the open edges joining its vertices.

\begin{definition}[Good clusters]\label{def:good}
Let $\varepsilon\in(0,1/2)$ and $\sigma\coloneqq1/2+\varepsilon$. A finite open cluster~$C$ with
$\abs{C}\le M$ is \emph{$M$-good} if the following holds for every $A\subseteq C$ and every connected
component $B$ of the open subgraph induced by $C\setminus A$: if the set $\mathcal K$ of open edges
between $B$ and $A$ is not empty, then $\diambr(B)\le\abs{\mathcal K}M^\sigma$.
\end{definition}

\begin{proposition}[Step~1a: bridges in critical clusters]\label{prop:good}
Assume \cref{ass:G} and let $\varepsilon\in(0,1/2)$. For all sufficiently large~$M$,
\begin{equation*}
  \PP_{p_c}\bigl(\abs{C_o}\le M \text{ and } C_o \text{ is not $M$-good}\bigr)\le\exp(-M^\varepsilon)\,.
\end{equation*}
Moreover, the event $\{\abs{C_o}\le M,\ C_o\text{ is $M$-good}\}$ is determined by the states of the
edges with an endpoint in $B(o,M)$.
\end{proposition}

In a critical branching process of total size~$M$, every edge is a bridge and the height is of order
$M^{1/2}$; heuristically, critical clusters on a nonamenable graph behave alike. \Cref{prop:good} is a
robust form of this heuristic: pieces cut out of a critical cluster have bridge diameter at most
$M^{1/2+\varepsilon}$ per cut edge, outside an event of stretched-exponentially small probability. It
is proved in \cref{sec:bridges} by an exploration-martingale bound on the number of pivotal edges of a
marked cluster. The exponent $\sigma$ enters the next step through the requirement $\beta+\sigma<1$.

\begin{definition}[Local endpoint maps]\label{def:map}
A \emph{local endpoint map} is a family $\Phi=(\Phi_x)_{x\in V}$ of functions
$\Phi_x\colon\{0,1\}^E\times V\to[0,1]$ with $\sum_y\Phi_x(\omega,y)\le1$, to be read as a randomised
rule: given the configuration~$\omega$, the rule at $x$ outputs $y$ with probability
$\Phi_x(\omega,y)$ and fails with the remaining probability. We require that there is $r_\Phi<\infty$
such that $\Phi_x(\omega,\cdot)$ depends only on the states of the edges with an endpoint in
$B(x,r_\Phi)$, that $\Phi_x(\omega,y)=0$ unless $y\in C_x(\omega)$, and that
$\Phi_{gx}(g\omega,gy)=\Phi_x(\omega,y)$ for all $g\in\Aut(G)$. The map has \emph{cap}~$M\in\N$ if
$\Phi_x(\omega,\cdot)\equiv0$ unless $\abs{C_x}\le M$ and $C_x$ is $M$-good. Its \emph{success
probability}, \emph{endpoint kernel} and \emph{kernel norm} are
\begin{equation*}
  s\coloneqq\EE_{p_c}\sum_{y}\Phi_x(\omega,y)\,,\qquad
  K(x,y)\coloneqq\frac{\EE_{p_c}\Phi_x(\omega,y)}{s}\,,\qquad
  r\coloneqq\norm{K}\,;
\end{equation*}
by invariance $s$ does not depend on~$x$. If $s>0$ and $r<1$, the \emph{efficiency ratio} of~$\Phi$ is
\begin{equation*}
  \operatorname{eff}(\Phi)\coloneqq\frac{\log(1/s)+\log M}{\log(1/r)}\,.
\end{equation*}
\end{definition}

Thus $K(x,y)$ is the probability that the rule at $x$ outputs~$y$, given that it succeeds; caps refer
to $M$-goodness for the parameter $\varepsilon$ of \cref{def:good}, which is kept fixed. Under
\cref{ass:U} the endpoint kernel is doubly stochastic, so $r\le1$ (\cref{lem:map-basic}). A small
efficiency ratio means that the selected vertex is spread out very widely, as measured on $\ell^2(V)$,
at a moderate price in success probability and cluster size. Steps of the lazy random walk along open
edges give a first map with $r<\e^{-1}$ (\cref{lem:start}); the point of the next proposition is to make
the efficiency ratio small.

\begin{proposition}[Step~1b: heavy tails yield efficient maps]\label{prop:amplification}
Assume \cref{ass:G,ass:U}, and suppose that $\EE_{p_c}\abs{C_o}^{\alpha}=\infty$ for some
$\alpha\in(0,1/2-\varepsilon)$. Then there are local endpoint maps with cap $M\ge2$ and $r<\e^{-1}$
whose efficiency ratios are arbitrarily small.
\end{proposition}

The proof in \cref{sec:maps} improves a given map~$\Phi$ in one \emph{amplification step}, at a heavy
scale~$n$, that is, a scale with $\PP_{p_c}(n\le\abs{C_o}<2n)\ge n^{-\beta}$ for some fixed
$\beta<1/2-\varepsilon$. For a root $x$, let $A$ be the cluster of $x$ in a first configuration, and
for a vertex $z$ run $\ell$ steps of $\Phi$ from~$z$, each in a fresh independent configuration. If
$n\le\abs A<2n$ and the last step lands in~$A$, we splice the pieces of the $\ell$ old clusters onto~$A$
by exchanging finitely many edge states between the configurations. Afterwards $z$ lies in the cluster
of~$x$, which has fewer than $2n+\ell M$ vertices. The new map at $x$ outputs $z$ with probability
proportional to the total weight of the splicings that produce the observed configuration and end
at~$z$. Its endpoint kernel is essentially the $\ell$-th power of the old one composed with a stochastic
kernel, so $\log(1/r)$ is multiplied by about~$\ell$. Since the old kernel is doubly stochastic, summing
over $z$ the probability that the chain from $z$ succeeds and ends in $A$ gives $s^\ell\abs A$; this
makes the new success probability of order $\PP_{p_c}(n\le\abs{C_o}<2n)\ge n^{-\beta}$, up to a factor
polynomial in $\ell$ and~$M$. The cap is about $2n$. To control fluctuations, the splicing must be
undone: once the positions of the $\ell$ seams among the bridges of the final cluster are known,
everything else is determined, and \cref{prop:good} bounds the number of possible seam positions,
weighted by the probabilities of choosing them, by about $M^{\sigma\ell}$. A second-moment estimate
needs this number to be at most $ns^\ell$, which is why the number of steps can be as large as about
$\log n/(\log(1/s)+\sigma\log M)$. With this choice, each amplification step multiplies the ratio
$(\log(1/s)+\sigma\log M)/\log(1/r)$ by a factor close to $\beta+\sigma<1$.

\begin{proposition}[Step~1c: efficient maps yield light tails]\label{prop:efficient}
Assume \cref{ass:G,ass:U}, and suppose that there are local endpoint maps with cap $M\ge2$ and
$r<\e^{-1}$ whose efficiency ratios are arbitrarily small. Then
$\EE_{p_c}\abs{C_o}^\alpha<\infty$ for every $\alpha\in(0,1/2)$.
\end{proposition}

The proof in \cref{sec:moments} runs a fixed efficient map at parameter $p$ slightly below $p_c$; its
cap is used only through $\abs{C_x}\le M$, not through goodness. Starting from the root of an
independent, size-biased cluster~$A$, iterates of the endpoint kernel spread so quickly that after
$\ell\approx\log\chi(p)/\log(1/r)$ steps the endpoint has left $A$, and the susceptibility in
$V\setminus A$ at the endpoint is comparable to $\chi(p)$. On the other hand, a successful output lies
within $M$ open steps of its input, and moving along at most $M$ open edges increases the
susceptibility in $V\setminus A$ by a factor at most of order $M$ (\cref{lem:intrinsic-ball}); after
dividing by the success probability, one step of the map multiplies it by at most a factor $Z\le CM/s$,
up to a contribution from the boundary of~$A$. The boundary contribution is controlled by the quantity
\begin{equation*}
  \Delta(p)\coloneqq\chi(p)^{-2}\sum_{b\sim o}\EE_p\bigl[\abs{C_o}\,\abs{C_b};\,C_o\neq C_b\bigr]
  =(1-p)\,\frac{\chi'(p)}{\chi(p)^2}\,,
\end{equation*}
where the identity holds for almost every $p$ by Russo's formula. Comparing the two bounds gives
$\Delta(p)\ge cZ^{-\ell}$, that is, $\chi'\ge c\chi^{2-\gamma}$ with $\gamma\approx\operatorname{eff}(\Phi)$,
and so $\chi(p)\le C\delta^{-1/(1-\gamma)}$. A comparison of explorations at $p$ and at $p_c$ then bounds
the critical tail: $\PP_{p_c}(\abs{C_o}\ge N)\le N^{-1/2+o(1)}$.

\begin{proof}[Proof of \cref{thm:moments} given \cref{prop:amplification,prop:efficient}]
Suppose that $\EE_{p_c}\abs{C_o}^{\alpha_0}=\infty$ for some $\alpha_0\in(0,1/2)$, and choose the
parameter of \cref{def:good} as some $\varepsilon\in(0,1/2-\alpha_0)$. By \cref{prop:amplification}, the
hypothesis of \cref{prop:efficient} holds, and \cref{prop:efficient} gives
$\EE_{p_c}\abs{C_o}^{\alpha_0}<\infty$, a contradiction.
\end{proof}

\begin{theorem}[Step~2: near-critical a priori bounds]\label{thm:apriori}
Assume \cref{ass:G,ass:U}. For every $\eta>0$ there is $C_\eta<\infty$ such that
\begin{align}
  \chi(p)&\le C_\eta\,(p_c-p)^{-1-\eta}&&\text{for all }p\in[0,p_c)\,,\label{eq:apriori-chi}\\
  \nabla_p&\le C_\eta\,(p_c-p)^{-\eta}&&\text{for all }p\in[p_c/2,p_c)\,.\label{eq:apriori-triangle}
\end{align}
\end{theorem}

Both bounds are proved from \cref{thm:moments}. The susceptibility bound~\eqref{eq:apriori-chi}
(\cref{sec:chi}) follows from a differential inequality for $\PP_p(\abs{C_o}\ge n)$ in~$p$, obtained from
the decision-tree inequality of O'Donnell, Saks, Schramm and Servedio, applied to an algorithm that
marks every vertex independently with probability $1/n$ and explores the clusters of the marked
vertices. The triangle bound~\eqref{eq:apriori-triangle} (\cref{sec:triangle}) follows from a
differential inequality for $\nabla_p$. By Russo's formula, $\nabla_p'$ is a sum over pairs of distinct
clusters, weighted by the number of edges between them and by a kernel built from $T_p$. Pairs with a
small cluster contribute at most a power of the size cutoff times $\nabla_p$. For two large clusters,
Step~1 bounds the probability that two adjacent distinct clusters are both large, and the number of
edges between them is controlled by comparing two explorations of one cluster, with and without the
other cluster removed. The centred scores of the two explorations differ by exactly $-p$ times this
number, and each score is concentrated.

\subsection{Part II: the corridor deformation (Steps 3 and 4)}\label{sec:part2}

We now define the deformation precisely. Throughout this subsection $p\in[p_c/2,p_c)$ and $k\ge3$ are
fixed, and $q\coloneqq1-p$.

\begin{definition}[Corridors and the deformed kernel]\label{def:deformation}
A \emph{raw path} is a sequence $w=(w_0,\dots,w_k)$ of vertices with $P(w_{i-1},w_i)>0$ for all~$i$;
it has \emph{base} $w_0$, \emph{end} $w_k$ and probability $\mu(w)\coloneqq\prod_{i=1}^kP(w_{i-1},w_i)$.
Delete the holding steps $w_{i-1}=w_i$ and erase the loops of the remaining path in chronological
order. This gives a self-avoiding path $\gamma_w$ from $w_0$ to $w_k$ with $\abs{\gamma_w}$ edges. Its
\emph{interior} $I_w$ is the set of its vertices other than $w_0$ and $w_k$, and $\Lambda_w\coloneqq
V\setminus I_w$. The \emph{blockers} $\Bl_w$ are the edges incident to $I_w$ that are not edges
of~$\gamma_w$. The raw path is a \emph{template} if $\abs{\gamma_w}\ge3$. Its \emph{corridor event}
$\Cor_w$ is the event that all edges of $\gamma_w$ are open and all blockers are closed, and its
\emph{weight} is
\begin{equation*}
  c(w)\coloneqq\mu(w)\,p^{-\abs{\gamma_w}}\,q^{-\abs{\Bl_w}}\,,\qquad\text{so that}\qquad
  c(w)\,\PP_p(\Cor_w)=\mu(w)\,.
\end{equation*}
For a domain $\Lambda\subseteq V$ and $x,y\in\Lambda$, a template
$w$ is \emph{counted} for $(x,y)$ in~$\Lambda$ if all vertices of $\gamma_w$ lie in~$\Lambda$, every
edge of $\gamma_w$ is open and \emph{pivotal} for $x\connin{\Lambda}y$ (that is, $x\connin{\Lambda}y$
holds and fails after closing this edge), and every blocker is closed or has an endpoint outside
$\Lambda$. With the sum over all templates~$w$ of any base, define
\begin{equation*}
  N^\Lambda_{xy}\coloneqq\sum_{w\text{ counted for }(x,y)\text{ in }\Lambda}c(w)\,,\qquad
  W^\Lambda_t(x,y)\coloneqq\1\{x\connin{\Lambda}y\}\,\e^{-tN^\Lambda_{xy}}\qquad(t\ge0)\,,
\end{equation*}
and $W^\Lambda_t(x,y)\coloneqq0$ if $x\notin\Lambda$ or $y\notin\Lambda$. The \emph{deformed
susceptibilities} and the \emph{deformed kernel} are
\begin{equation*}
  \chi^\Lambda_t(x)\coloneqq\EE_p\sum_{y\in V}W^\Lambda_t(x,y)\,,\qquad
  \chi_t\coloneqq\chi^V_t(o)\,,\qquad
  T_t(x,y)\coloneqq\EE_p\,W^V_t(x,y)\,.
\end{equation*}
Finally, $\Erw[F(w)]\coloneqq\sum_w\mu(w)F(w)$ denotes the average over $k$-step lazy-walk paths $w$ with
base $w_0=o$.
\end{definition}

At $t=0$ nothing is penalised: $T_0=T_p$ and $\chi_0=\chi(p)$. Different raw paths with the same
loop erasure are different templates, and the convention for blockers outside~$\Lambda$ lets us compare
domains without changing the weights.

\begin{figure}[ht]
\centering
% A corridor on the event Cor_w (schematic)
\begin{tikzpicture}[scale=0.95,
    vtx/.style={circle, fill, inner sep=1.3pt},
    blob/.style={draw=figgrey, fill=figgrey!12, rounded corners=14pt},
    corr/.style={figmain, line width=1.6pt},
    blk/.style={figgrey, dashed, thick},
    leg/.style={black, thick, decorate, decoration={snake, amplitude=0.9pt, segment length=6pt}}]
  % exterior clusters
  \draw[blob] (-2.6,-1.25) rectangle (1.5,1.25);
  \draw[blob] (8.1,-1.25) rectangle (12.2,1.25);
  \node[font=\small] at (-0.55,-1.6) {$C_u^{\Lambda_w}$};
  \node[font=\small] at (10.15,-1.6) {$C_v^{\Lambda_w}$};
  % corridor
  \coordinate (u) at (1.1,0);
  \coordinate (a) at (2.9,0.55);
  \coordinate (b) at (4.8,0.75);
  \coordinate (c) at (6.7,0.45);
  \coordinate (v) at (8.5,0);
  \draw[corr] (u) -- (a) -- (b) -- (c) -- (v);
  % closed blockers at interior vertices
  \foreach \p/\ang in {a/100, a/250, a/300, b/80, b/120, b/265, c/70, c/240, c/290} {
    \draw[blk] (\p) -- ++(\ang:0.6);
  }
  \foreach \p in {a,b,c} \node[vtx, figmain] at (\p) {};
  \node[vtx] at (u) {};
  \node[vtx] at (v) {};
  \node[font=\small, anchor=north east] at (1.15,-0.05) {$u$};
  \node[font=\small, anchor=north west] at (8.45,-0.05) {$v$};
  \node[font=\small, text=figmain] at (4.8,1.75) {interior $I_w$};
  % terminals and exterior legs
  \coordinate (x) at (-1.9,0.35);
  \coordinate (y) at (11.5,-0.35);
  \draw[leg] (x) -- (-0.4,-0.45) -- (u);
  \draw[leg] (v) -- (9.9,0.45) -- (y);
  \node[vtx] at (x) {};
  \node[vtx] at (y) {};
  \node[font=\small, anchor=south] at (x) {$x$};
  \node[font=\small, anchor=south] at (y) {$y$};
\end{tikzpicture}

\caption{A corridor on the event $\Cor_w$ (schematic). The loop-erased path $\gamma_w$ (thick) is
open, and every other edge at its interior $I_w$ is closed (dashed). Its edges are all pivotal for
$x\conn y$ if and only if $x$ and $y$ lie in the clusters $C_u^{\Lambda_w}$ and $C_v^{\Lambda_w}$ of its
endpoints in $\Lambda_w=V\setminus I_w$, in either order, and these two clusters are distinct
(\cref{lem:calibration}).}\label{fig:corridor}
\end{figure}

\begin{proposition}[Step~3: properties of the deformation]\label{prop:deformation}
Assume \cref{ass:G,ass:U}. There is $C_0\ge1$, depending only on~$G$, such that, with
$c_k\coloneqq\e^{C_0k}$, the following holds for all $p\in[p_c/2,p_c)$, $k\ge3$ and $t\ge0$.
\begin{enumerate}[label=\textup{(\alph*)}]
  \item\label{it:def-basic} \emph{Basic bounds.} $T_t$ is symmetric and invariant,
  $0\le T_t\le T_p$ entrywise, and $\norm{T_t}\le\chi_t\le\chi(p)$. The function $t\mapsto\chi_t$ is
  continuously differentiable on $[0,\infty)$.
  \item\label{it:def-norm} \emph{The norm barely moves.}
  $\norm{T_t}\ge\norm{T_p}-2t\rho^k\norm{T_p}^2$.
  \item\label{it:def-rate} \emph{The row sum decreases.}
  \begin{equation*}
    -\frac{\d}{\d t}\chi_t\ge\e^{-tc_k}\Bigl(\Erw\bigl[\1\{\abs{\gamma_w}\ge3\}\,
    \chi^{\Lambda_w}_t(w_0)\,\chi^{\Lambda_w}_t(w_k)\bigr]-C_0\,\chi(p)^2\rho^k\,\nabla_p\Bigr)\,.
  \end{equation*}
  \item\label{it:def-deletion} \emph{Deletion.} For every domain $\Lambda$, every $A\subseteq V$ and
  every $z\in V$,
  $\chi^{\Lambda\setminus A}_t(z)\ge\chi^\Lambda_t(z)-\chi(p)\sum_{b\in\Nb[A]}T_p(z,b)$.
  \item\label{it:def-path} \emph{Transport along a path.} If $\Lambda$ contains a self-avoiding path
  of length~$l$ from $b$ to~$z$, then $\chi^\Lambda_t(b)\ge\e^{-tlc_k}p^l\chi^\Lambda_t(z)$.
\end{enumerate}
\end{proposition}

Items~\ref{it:def-norm} and~\ref{it:def-rate} are the two halves of the comparison described in
\cref{sec:mechanism}. Item~\ref{it:def-norm} holds because a penalised connection factors through a
corridor, so that $T_p-T_t\le2t\,T_pP^kT_p$ entrywise. Item~\ref{it:def-rate} is obtained by moving
the root of the susceptibility to the base of the corridor with the mass-transport principle: the
calibration $c(w)\PP_p(\Cor_w)=\mu(w)$ then produces the average over raw paths, and the two exterior
clusters contribute the product of the deformed susceptibilities at the two ends. The error term
accounts for the requirement that the exterior clusters be distinct and is bounded with the
van den Berg--Kesten inequality and the trace. The factor $\e^{-tc_k}$ is the price for the penalties
of corridors that overlap the given one. Items~\ref{it:def-deletion} and~\ref{it:def-path} are the
tools that make the main term of~\ref{it:def-rate} usable: deleting a set $A$ can only destroy the
connections that pass near~$A$, and an open path transports mass at the cost of its probability.
\Cref{prop:deformation} is proved in \cref{sec:deformation}.

\begin{proposition}[Step~4: escapes at both ends]\label{prop:escapes}
Assume \cref{ass:G,ass:U}. There are $J,k_0\in\N$ and $C_1,c_\star>0$, depending only on~$G$, with the
following property. Let $p\in[p_c/2,p_c)$, $k\ge k_0$ and $R\ge1$, and put
$j\coloneqq J+\lceil C_1\log((k+1)R)\rceil$. Let $\mathcal W$ be the set of raw paths $w$ with
$\abs{\gamma_w}\ge3$ such that each of $w_0$ and $w_k$ is joined by a self-avoiding path in
$\Lambda_w$ of length at most $j$ to a vertex~$z$ with
\begin{equation*}
  \sum_{b\in\Nb[I_w]}T_p(z,b)\le\frac{1}{4R}\,.
\end{equation*}
Then $\Erw[\1\{w\in\mathcal W\}]\ge c_\star$.
\end{proposition}

The vertices $z$ are \emph{escapes}: from them, deleting the corridor costs little in
\cref{prop:deformation}\ref{it:def-deletion}. The proof in \cref{sec:escapes} lets the raw path begin
and end with halves of geodesic segments, and lets the escapes leave along the opposite halves and then
continue as independent random walks of length $m$ of order $C_1\log((k+1)R)$. Avoidance of the raw
path follows from the decay of $P^l$ on $\ell^2(V)$. The construction is symmetric under reversal of the
raw path, and the mass-transport principle turns this symmetry into an identity of probabilities, which
transfers the avoidance estimate from one endpoint to the other. Finally, the walks smooth the
connection kernel, $(P^mT_p)(a,b)\le\rho^m$, which makes the sum over $\Nb[I_w]$ small.

\subsection{Proof of critical boundedness}\label{sec:assembly}

We now prove \cref{thm:crit-bdd} from the four steps and the sprinkling bound.

\begin{proof}[Proof of \cref{thm:crit-bdd}]
Suppose, for a contradiction, that $\norm{T_{p_c}}=\infty$. For $p\in[p_c/2,p_c)$ put
\begin{equation*}
  \delta\coloneqq p_c-p\,,\qquad S\coloneqq\log(1/\delta)\,,\qquad R\coloneqq\frac{\chi(p)}{\norm{T_p}}\ge1\,.
\end{equation*}
All limits below are taken as $\delta\downarrow0$, that is, $p\uparrow p_c$.

\emph{Stage 1: size of the parameters.}
By~\eqref{eq:norm-lower}, which follows from \cref{lem:sprinkling}, $\norm{T_p}\ge(1-p_c)/(d\delta)$,
so $\log\norm{T_p}\ge S-C$. By
\cref{thm:apriori}, for every $\eta>0$ we have $R\le C_\eta\delta^{-1-\eta}\cdot d\delta/(1-p_c)\le
C'_\eta\delta^{-\eta}$ and $\nabla_p\le C_\eta\delta^{-\eta}$. Hence
\begin{equation}\label{eq:assembly-subpower}
  \log R=o(S)\qquad\text{and}\qquad\log\nabla_p=o(S)\,.
\end{equation}

\emph{Stage 2: choice of the constants.}
Fix the constants $J,k_0,C_1,c_\star$ of \cref{prop:escapes}; they do not depend on~$p$. Since $p\ge p_c/2$ and $j\le J+1+C_1\log((k+1)R)$, we have
\begin{equation}\label{eq:assembly-pj}
  p^{2j}\ge c_2\,(kR)^{-C_2}\qquad\text{for all }k\ge1\text{ and }R\ge1\,,
\end{equation}
with $C_2\coloneqq2C_1\log(2/p_c)$ and $c_2\coloneqq(p_c/2)^{2J+2}\,2^{-C_2}$. Let
$\lambda\coloneqq\log(1/\rho)>0$, fix $K_\star>0$ with $\kappa_\star\coloneqq\lambda K_\star>C_2+4$, and set
\begin{equation*}
  A_\delta\coloneqq\log(2R)+\log(2+\nabla_p)+\log S\,,\qquad
  k\coloneqq\lceil K_\star A_\delta\rceil\,,\qquad
  D_\star\coloneqq(kR)^{C_2+2}\,,\qquad
  t_\star\coloneqq\frac{D_\star}{\norm{T_p}}\,.
\end{equation*}
By~\eqref{eq:assembly-subpower}, $k\to\infty$ and $k=o(S)$; in particular $k\ge\max\{k_0,3\}$ and
$k\le S$ for small~$\delta$. We claim that
\begin{equation}\label{eq:assembly-small}
  D_\star\rho^k\to0\,,\qquad(kR)^{C_2}R^2\nabla_p\,\rho^k\to0\,,\qquad t_\star jc_k\to0\,.
\end{equation}
Indeed, $\rho^k\le\e^{-\lambda K_\star A_\delta}=(2R)^{-\kappa_\star}(2+\nabla_p)^{-\kappa_\star}
S^{-\kappa_\star}$. Since $R,\nabla_p\ge1$ and $k\le S$, this gives
$D_\star\rho^k\le S^{C_2+2-\kappa_\star}\to0$ and
$(kR)^{C_2}R^2\nabla_p\,\rho^k\le S^{C_2-\kappa_\star}\to0$. Finally, $c_k=\e^{C_0k}$ and
$\log\norm{T_p}\ge S-C$ give
\begin{equation*}
  \log(t_\star jc_k)\le(C_2+2)(\log k+\log R)-S+C+\log j+C_0k=-S+o(S)\,,
\end{equation*}
by~\eqref{eq:assembly-subpower}, $k=o(S)$ and $j=O(\log(kR))$.

\emph{Stage 3: the norm stays large.}
Consider the deformation of \cref{def:deformation} with these $p$ and~$k$, for $t\in[0,t_\star]$. By
\cref{prop:deformation}\ref{it:def-norm} and~\eqref{eq:assembly-small}, for small $\delta$,
\begin{equation}\label{eq:assembly-norm}
  \norm{T_t}\ge\norm{T_p}\bigl(1-2t_\star\rho^k\norm{T_p}\bigr)=\norm{T_p}(1-2D_\star\rho^k)\ge\tfrac12\norm{T_p}\,,
  \qquad\text{hence}\qquad
  \chi_t\ge\tfrac12\norm{T_p}
\end{equation}
by \cref{prop:deformation}\ref{it:def-basic}.

\emph{Stage 4: the row sum decreases quadratically.}
Let $w\in\mathcal W$, the set of \cref{prop:escapes} for our $p$, $k$ and~$R$, and let $z$ be the
escape vertex of the endpoint $w_0$. The deformation is invariant, so $\chi_t^V(z)=\chi_t$. By
\cref{prop:deformation}\ref{it:def-deletion} with $\Lambda=V$ and $A=I_w$, then $\chi(p)=R\norm{T_p}$
and~\eqref{eq:assembly-norm},
\begin{equation*}
  \chi^{\Lambda_w}_t(z)\ge\chi_t-\frac{\chi(p)}{4R}=\chi_t-\tfrac14\norm{T_p}\ge\tfrac12\chi_t\,.
\end{equation*}
\Cref{prop:deformation}\ref{it:def-path}, applied along the escape path of length at most~$j$ in
$\Lambda_w$, gives $\chi^{\Lambda_w}_t(w_0)\ge\e^{-tjc_k}p^j\chi_t/2$, and the same bound holds at~$w_k$.
Since $\e^{-2t_\star jc_k}\ge1/2$ for small~$\delta$ by~\eqref{eq:assembly-small}, the bound
\eqref{eq:assembly-pj} and $\Erw[\1\{w\in\mathcal W\}]\ge c_\star$ yield, for all $t\in[0,t_\star]$,
\begin{equation*}
  \Erw\bigl[\1\{\abs{\gamma_w}\ge3\}\,\chi^{\Lambda_w}_t(w_0)\,\chi^{\Lambda_w}_t(w_k)\bigr]
  \ge c_\star\cdot\tfrac12\,p^{2j}\cdot\tfrac14\chi_t^2
  \ge\tfrac18c_\star c_2\,(kR)^{-C_2}\chi_t^2\,.
\end{equation*}
The error term in \cref{prop:deformation}\ref{it:def-rate} satisfies, by $\chi(p)=R\norm{T_p}\le2R\chi_t$
and~\eqref{eq:assembly-small},
\begin{equation*}
  C_0\,\chi(p)^2\rho^k\nabla_p\le4C_0\,R^2\nabla_p\,\rho^k\,\chi_t^2=o\bigl((kR)^{-C_2}\bigr)\,\chi_t^2\,.
\end{equation*}
Since also $\e^{-tc_k}\ge1/2$ for $t\le t_\star$, \cref{prop:deformation}\ref{it:def-rate} gives, for
small~$\delta$,
\begin{equation}\label{eq:assembly-ode}
  -\frac{\d}{\d t}\chi_t\ge c_3\,(kR)^{-C_2}\chi_t^2\qquad\text{for all }t\in[0,t_\star]\,,
  \qquad c_3\coloneqq c_\star c_2/32\,.
\end{equation}

\emph{Stage 5: contradiction.}
The function $t\mapsto\chi_t$ is positive and continuously differentiable, so~\eqref{eq:assembly-ode}
integrates to
\begin{equation*}
  \frac{1}{\chi_{t_\star}}\ge\frac1{\chi_0}+c_3\,(kR)^{-C_2}t_\star\ge c_3\,\frac{(kR)^2}{\norm{T_p}}\,.
\end{equation*}
By~\eqref{eq:assembly-norm}, $1/\chi_{t_\star}\le2/\norm{T_p}$. Hence $c_3(kR)^2\le2$, which fails for
small~$\delta$ because $k\to\infty$. This contradiction proves $\norm{T_{p_c}}<\infty$.
\end{proof}

\begin{remark}[Order of the choices]\label{rem:order}
The constants $J,k_0,C_1,c_\star$ of the escape construction and $C_0$ of the deformation depend only
on~$G$; they are fixed before $p$ is chosen. The length $k$ and the final time $t_\star$ are then chosen
as functions of $p$ through $R$, $\nabla_p$ and $S$. The a priori bounds of Step~2 are used only
through~\eqref{eq:assembly-subpower}.
\end{remark}

\subsection{Roadmap}\label{sec:roadmap}

\Cref{fig:dependencies} shows how the statements depend on each other. Part~I
(\cref{sec:bridges,sec:maps,sec:moments,sec:triangle}) establishes the a priori estimates. It only
concerns critical and near-critical percolation and does not involve the deformation. Part~II
(\cref{sec:deformation,sec:escapes}) holds at every fixed $p\in[p_c/2,p_c)$ and uses only subcritical
facts; the two parts meet only in the proof above.
